\section{Diagram Optimization via Categorical Rewriting}

Cells do not merely execute metabolic pathways---they \emph{optimize} them.
Flux Balance Analysis identifies rate-limiting enzymes, allosteric regulation
redirects flux through cheaper branches, and pathway rewiring eliminates
dead-end metabolites. The result is not a different pathway, but a more
efficient execution of the same input-output behavior. In this section we
show that wiring diagrams admit analogous optimizations: cost annotations
make resource consumption explicit, static analysis identifies structural
opportunities, and rewriting rules transform diagrams while preserving
observational equivalence.

\subsection{Cost-Annotated Diagrams}

We extend the WAgent operad (Section~4) with resource annotations.
Define a \emph{resource cost} monoid $(\mathbb{R}_{\geq 0}^3, +, \mathbf{0})$
where each element is a triple $(\text{atp}, \text{latency}, \text{memory})$.
A cost-annotated wiring diagram enriches the category of wiring diagrams
over this monoid: each module $m$ carries a cost $c(m) \in \mathbb{R}_{\geq 0}^3$,
and each wire $w$ carries a transmission cost $c(w) \in \mathbb{Z}_{\geq 0}$
(measured in ATP units).

\begin{definition}[Cost-Annotated Wiring Diagram]
A \emph{cost-annotated wiring diagram} is a tuple $(M, W, c_M, c_W)$ where:
\begin{itemize}[leftmargin=*]
\item $M$ is a finite set of modules, each with typed input/output ports;
\item $W \subseteq \{(m_s.p_s, m_d.p_d) \mid m_s, m_d \in M\}$ is a set of wires;
\item $c_M : M \to \mathbb{R}_{\geq 0}^3$ assigns resource costs to modules;
\item $c_W : W \to \mathbb{Z}_{\geq 0}$ assigns transmission costs to wires.
\end{itemize}
The \emph{path cost} of a directed path $m_1 \xrightarrow{w_1} m_2 \xrightarrow{w_2} \cdots \xrightarrow{w_{n-1}} m_n$
is $\sum_{i=1}^{n} c_M(m_i) + \sum_{i=1}^{n-1} c_W(w_i)$.
\end{definition}

Biologically, $c_M$ corresponds to the ATP cost of enzyme catalysis,
while $c_W$ models the energetic cost of metabolite transport between
cellular compartments.

\subsection{Rewriting Rules as Endofunctors}

Each optimization pass is an endofunctor $F : \mathbf{WDiag} \to \mathbf{WDiag}$
on the category of wiring diagrams that preserves input-output behavior.
Two diagrams $D_1, D_2$ are \emph{behaviorally equivalent} if for all
input assignments they produce the same outputs---i.e., they are bisimilar
in the coalgebraic sense of Section~3.5.

\begin{theorem}[Optimization Soundness]
Let $F$ be an optimization pass. If $F$ only removes wires that provably
never transmit data (dead wire elimination) or reorders modules within
the same topological layer (cost-order scheduling), then $F(D)$ is
bisimilar to $D$ for all diagrams $D$.
\end{theorem}

\noindent Three concrete passes are implemented:

\begin{enumerate}[leftmargin=*]
\item \textbf{Dead Wire Elimination.} A wire carrying a PrismOptic
whose accepted types have no intersection with the source port's
DataType can never transmit. Removing it does not change behavior.
This corresponds to pruning vestigial metabolic branches---enzymes
that never encounter their substrate are downregulated.

\item \textbf{Parallel Grouping.} Modules with no mutual dependencies
form a \emph{parallel group} and can execute concurrently. This is
identified via topological layering of the dependency DAG. Biologically,
independent metabolic pathways (e.g., glycolysis and fatty acid oxidation)
operate simultaneously in different cellular compartments.

\item \textbf{Cost-Order Scheduling.} Within each parallel group,
modules are sorted by ascending cost. Cheaper modules execute first,
allowing early termination or budget reallocation if expensive modules
would exceed available ATP. This mirrors cellular enzyme kinetics:
low-$K_m$ (high-affinity, low-cost) enzymes are preferentially activated.
\end{enumerate}

\subsection{Critical Path Analysis}

The \emph{critical path} of a cost-annotated diagram is the longest
cost-weighted path through the dependency DAG:

\[
\text{CP}(D) = \arg\max_{P \in \text{Paths}(D)} \sum_{m \in P} c_M(m) + \sum_{w \in P} c_W(w)
\]

The critical path determines the minimum execution time under maximal
parallelism---no schedule can complete faster than the critical path
cost. In metabolic terms, this is the rate-limiting pathway: the
bottleneck whose throughput constrains the entire system. Identifying
it enables targeted optimization (e.g., caching expensive modules,
splitting them into cheaper sub-diagrams).

\subsection{Resource-Aware Execution}

The \texttt{ResourceAwareExecutor} gates module execution on MetabolicState,
implementing the biological principle that pathway activity is regulated
by cellular energy availability:

\begin{itemize}[leftmargin=*]
\item \textbf{STARVING}: Only essential modules execute; non-essential
modules are skipped. This mirrors the starvation response where cells
shut down biosynthetic pathways to conserve ATP for survival functions.

\item \textbf{CONSERVING}: Expensive non-essential modules are deferred.
Cheap modules execute normally. This corresponds to metabolic triage
under moderate stress.

\item \textbf{NORMAL/FEASTING}: Full execution with parallel scheduling.
Independent module groups execute concurrently via thread pools. This
mirrors the fed state where abundant ATP enables all pathways.
\end{itemize}

\noindent A \texttt{BudgetOptic} provides wire-level cost capping:
once cumulative transmission cost through a wire exceeds a threshold,
further data flow is blocked. This implements pathway-level budget
caps analogous to allosteric feedback inhibition.

\subsection{Relation to Abbott \& Zardini}

Abbott and Zardini~\cite{abbott2025categorical} demonstrate that
category-theoretic string diagrams are optimizable formal objects,
deriving FlashAttention ``on a napkin'' through diagram analysis.
Our approach extends this insight in two directions:

\begin{enumerate}[leftmargin=*]
\item \textbf{Dynamic state.} Where Abbott \& Zardini optimize static
algorithm pipelines, operon optimizes \emph{coalgebraic} systems
with internal state (Section~3.5). The MetabolicState-gated executor
makes optimization decisions that depend on runtime state---a feature
absent from purely structural diagram analysis.

\item \textbf{Resource enrichment.} Our cost annotations form a
monoidal enrichment of the wiring category, making resource consumption
a first-class categorical concept rather than an external analysis.
This enables compositional cost reasoning: the cost of a composite
diagram is determined by the costs of its components.
\end{enumerate}

\noindent Both approaches share the core insight that diagram structure
reveals optimization opportunities invisible to sequential code analysis.
The categorical framework ensures that optimizations preserve the
compositional semantics of the system.
